\documentclass[10pt,a4paper]{amsart}
\usepackage[margin=19mm]{geometry}
\usepackage{amsmath,amssymb,mathtools}
\usepackage[scr=rsfs,cal=euler]{mathalfa}
\usepackage{xfrac}
\usepackage[unicode,hidelinks,bookmarks=false]{hyperref}
\newcommand{\ie}{i.e.\ }
\newcommand{\cf}{cf.\ }
\newcommand{\resp}{respectively\ }
\newcommand{\Subsection}{\subsection}
\providecommand{\nobreakdash}{\nobreak\hbox{-}}
\setlength{\emergencystretch}{2em}
\multlinegap0pt
\usepackage{bm}
\usepackage{bbm}
\usepackage{dsfont}
\usepackage{xspace}
\usepackage{cancel}
\usepackage{mathtools}
\usepackage{cleveref}
\usepackage{amsthm}
\makeatletter
\@ifundefined{theorem}{\newtheorem{theorem}{Theorem}}{}
\@ifundefined{lemma}{\newtheorem{lemma}[theorem]{Lemma}}{}
\@ifundefined{proposition}{\newtheorem{proposition}[theorem]{Proposition}}{}
\makeatother
\DeclareMathOperator{\len}{\ell}
\newcommand{\wt}[1]{\widetilde{ #1 }}
\title[The geometry of subgroups of mapping tori of free groups]{The geometry of subgroups of mapping tori of free groups}
\author{Marco Linton}
\date{Source: arXiv:2510.03145v1 (CC BY 4.0)}
\begin{document}
\maketitle

\noindent\emph{Source and licence.} The geometry of subgroups of mapping tori of free groups. Version arXiv:2510.03145v1 (CC BY 4.0), distributed under the Creative Commons Attribution 4.0 International licence, as recorded in the arXiv licence metadata for this exact version and shown on the abstract page https://arxiv.org/abs/2510.03145v1. Full rights notice: licenses/arxiv-2510.03145-CC-BY-4.0.txt.\par\smallskip

\noindent\textit{Editorial scope.} Counted unit: the Proposition whose source label is \texttt{prop:hallways} in the article (source0.tex lines 1351--1358, source label \texttt{prop:hallways}), delivered together with the article's own complete original proof of it (source0.tex lines 1360--1432, one \texttt{proof} environment of 73 lines). No mathematics is written, repaired or re-derived: statement and proof are the article's own text line for line. Required context retained uncounted: lem:iteration\_general (lines 1301--1306); lem:iteration (lines 1339--1342). Every definition, display and label of those lines is retained unchanged. Omitted from this package and not redistributed: 1--1300, 1307--1338, 1343--1350, 1359--1359, 1433--1917. Nothing inside a retained excerpt is omitted or altered: every retained body is delivered exactly as the article's lines are written, so no omission receipt of any kind accompanies this package. Citations: the retained excerpts carry no citation command at all, so no bibliography entry of the article is transcribed and no reference is redistributed. Reference adaptation: none was needed and none was made; every reference inside the delivered unit resolves inside it, so no unresolved reference or citation is produced. Printed slips kept literal: no printed word, symbol, hypothesis or number of the counted statement or of its proof was altered. Layout and class: the article's own class is \texttt{amsart} with packages including amssymb, amsmath, amsthm, tikz-cd, graphicx, hyperref, mathtools, appendix, caption, subcaption, blindtext, babel, mathrsfs, mathabx; the wrapper is the lane's pinned \texttt{amsart} editorial wrapper at 10pt on A4, which restates the theorem-like environments the retained text needs and adds the article's own macro definitions that the retained text needs (\texttt{\textbackslash len}, \texttt{\textbackslash wt}), with extra wrapper packages (bm, bbm, dsfont, xspace, cancel, mathtools, cleveref). The delivered numbering is the unit's own. Assets: no retained span contains a figure environment, a graphics-inclusion command, a PDF-page inclusion, a table or a TikZ picture; no image file is copied into this package, the package declares no asset dependency and no image file is redistributed with it. Licence: version arXiv:2510.03145v1 of this article is distributed under the Creative Commons Attribution 4.0 International licence, as recorded in the arXiv licence metadata for this exact version and shown on the abstract page https://arxiv.org/abs/2510.03145v1. A full rights notice is delivered at licenses/arxiv-2510.03145-CC-BY-4.0.txt. Characters: every retained excerpt is pure ASCII, so the delivered engine, XeLaTeX with its default Latin Modern text font, reports no missing character.
\begin{lemma}
\label{lem:iteration_general}
Let $\omega\geqslant 0$ and let $m\geqslant k$ be an integer. There is a constant $M$ such that the following holds. 

Let $\gamma_1, \gamma_2\colon I \to R$ be immersed paths with $\len(\gamma_1), \len(\gamma_2)\leqslant \omega$ and let $\lambda\colon I\to R_{l}$ be an immersed path such that $\gamma_1*f^m(\lambda)*\gamma_2$ is path homotopic to an immersed path $\delta\colon I\to R_{l}$. If $\len(\delta)\geqslant M$, then $\delta = \delta_1*\alpha*\delta_2$ with $\len(\delta_1), \len(\delta_2)\leqslant M$ and with $\alpha$ a non-trivial path supported in $\bigsqcup_{i=1}^nR_{A_i}$.
\end{lemma}

\begin{lemma}
\label{lem:iteration}
Let $m\geqslant k$ be an integer and let $\lambda\colon I\to R_l$ be a path connecting $v_i$ with $v_j$ for some $1\leqslant i, j\leqslant n$. If $f^m(\lambda)$ is path homotopic to an immersed path $\delta\colon I \to R_l$, then $\sigma^m(i) = \sigma^m(j)$ and $\delta$ is supported in $R_{A_{\sigma^m(i)}}$. 
\end{lemma}

\begin{proposition}
\label{prop:hallways}
If $m\geqslant k$ and $\rho\geqslant 0$, then there is a constant $L\geqslant 0$ such that the following holds. If $h\colon [-m, m]\times I\to \overline{M}_{l}$ is an essential hallway, then
\begin{itemize}
\item $h$ is not cone-bounded,
\item the girth of $h$ is at most $L$ if $h$ is $\rho$-thin.
\end{itemize}
\end{proposition}

\begin{proof}
After performing a homotopy to $h$, we may assume that the path $h\mid[-m, m]\times \{0\}$ is of the form
\[
t_{-m}*g_{-m+1}*\ldots*g_0*t_0*g_1*t_1*\ldots*g_m*t_m
\]
where each $g_i$ is supported in a copy of $\wt{R}_{l}$, where $t_{-m}$ is a half edge of the form $[0, 1]\times\{x_{-m}\}$, $t_m$ is a half edge of the form $[-1, 0]\times\{x_m\}$ and where $t_i$ is an edge of the form $[-1, 1]\times\{x_i\}$ for all other values of $i$ and where here $x_i\in X_{e_i}^{(0)}$ is a 0-cell in an edge space associated with the edge $\pi(t_i) = e_i\in E(T)$. Note that this homotopy only increases lengths of the paths $h\mid[i, i+1]\times\{0\}$ by a constant $\xi$ depending on the metrics on the 2-cells in $\overline{M}_l$ (of which there are finitely many types). Similarly, we may assume that $h\mid[0, m]\times \{1\}$ is a path of the form
\[
t'_{-m}*g'_{-m+1}*\ldots*g'_0*t'_0*g'_1*t'_1*\ldots*g'_m*t'_m.
\]
If $h$ was $\rho$-thin before the homotopy, $h$ will by $\xi\rho$-thin after the homotopy. For ease of notation, replace $\rho$ with $\xi\rho$.

For each $i$, denote by $h_i^{\pm}$ the geodesic in $X_{e_i^{\pm}}$ connecting the endpoints of the path $\partial^{\pm}_{e_i}\circ h\mid\{i\}\times I$.

Suppose first for a contradiction that $h$ is cone-bounded. Then each $g_i, g_i'$ is trivial by definition and so $f^m(\mathfrak{p}\circ h_0^+)$ is path homotopic to $\mathfrak{p}\circ h_m^+$. Now \cref{lem:iteration} implies that $h_m^+$ is supported in a copy of the coning off of the universal cover of $R_{A_i}$ for some $i$. This implies that $h\mid\{m\}\times I$ is a geodesic connecting a cone point with itself which is not possible (since it would have to be trivial), a contradiction. Thus, $h$ cannot be cone-bounded as claimed.

Now suppose that $h$ is $\rho$-thin. Then each $g_{i}$, $g_{i}'$ has length at most $\rho$. By definition of $K, C$, we have that 
\begin{equation}
\label{eq: length bound 1}
K^{-1}\cdot \len(h\mid\{i\}\times I)  - C\leqslant \len(h_i^{\pm})\leqslant K\cdot \len(h\mid\{i\}\times I) + C
\end{equation}
for each $i$. The map $h|[i, i+1]\times I$ implies that 
\begin{equation}
\label{eq: path_homotopy}
g_{i+1}*h_i^+*\overline{g}_{i+1}'\sim h_{i+1}^-.
\end{equation}
where $\sim$ denotes path homotopy within the corresponding vertex space. We have
\begin{align*}
\len(h_i^+) - 2\rho \leqslant \len(h_i^+) - \len(g_i) - \len(g_i') &\leqslant \len(h_{i+1}^-)\\
						&\leqslant \len(h_i^+) + \len(g_i) + \len(g_i') \leqslant \len(h_i^+) + 2\rho. 
\end{align*}
Combining this with \eqref{eq: length bound 1} we see that
\begin{align*}
\len(h\mid\{i\}\times I)&\leqslant K\len(h^{\pm}_i) + KC\\
					&\leqslant K(\len(h^{\mp}_{i\pm1}) + 2\rho) + KC\\
					&\leqslant K(K\len(h\mid\{i\pm1\}\times I) + C + 2\rho) + KC\\
					&\leqslant K^2\len(h\mid\{i\pm1\}\times I) + 2K(C + \rho).
\end{align*}
Thus, by induction on $-m\leqslant n\leqslant m$, we have
\begin{align}
\label{eq: length bound 2}
K^{-2|n|}\cdot \len(h\mid\{0\}\times I) - 2(C+\rho)&\leqslant \len(h\mid\{n\}\times I)\\
\label{eq: length bound 3}				&\leqslant K^{2|n|}\cdot(\len(h\mid\{0\}\times I) + 2(C+\rho)).
\end{align}
Now suppose that for some $-m\leqslant i< m$, the paths $\overline{t}_i$ and $t_{i+1}$ follow the induced orientation on $T$. Then $\partial_{e_i}^+(X_{e_i})$ and $\partial_{e_{i+1}}^-(X_{e_{i+1}})$ both project to $R_{l-1}$ under the cover $c$. Since $R_{l-1}$ is a subgraph of $R_{l}$, we see that either $\partial_{e_i}^+(X_{e_i}) = \partial_{e_{i+1}}^-(X_{e_{i+1}})$ and $e_i = \overline{e}_{i+1}$ or $\partial_{e_i}^+(X_{e_i})\cap\partial_{e_{i+1}}^-(X_{e_{i+1}}) = \emptyset$ and $e_i\neq \overline{e}_{i+1}$. But since $h$ is an essential hallway, we must be in the latter case. The paths $g_{i+1}$ and $g_{i+1}'$ thus connect the two disjoint subtrees $\partial_{e_i}^+(X_{e_i})$ and $\partial_{e_{i+1}}^-(X_{e_{i+1}})$. Hence $g_{i+1}, g_{i+1}'$ exit $\partial_{e_i}^+(X_{e_i})$ at the same point and enter $\partial_{e_{i+1}}^-(X_{e_{i+1}})$ at the same point. This implies that $\len(h_{i+1}^-), \len(h_i^+)\leqslant 2\rho$. Combining with \eqref{eq: length bound 1}, we see that $\len(h\mid \{i\}\times I)\leqslant K(2\rho +C)$. By \eqref{eq: length bound 2} we see that
\[
\len(h\mid \{0\}\times I)\leqslant 2(K^{2m}\rho + C + \rho)
\]

Now assume that no such $i$ exists. In particular, there is at most one $-m\leqslant i< m$ so that $t_i$ and $\overline{t}_{i+1}$ follow the induced orientation on $T$ and for all other $i$, either $t_i, t_{i+1}$ or $\overline{t}_i, \overline{t}_{i+1}$ follow the induced orientation on $T$. In any case, there is an $i$ so that (after possibly flipping $h$), the edges $t_i, t_{i+1}, \ldots, t_{i+m}$ all follow the induced orientation on $T$.

The null-homotopy $h\mid[i, i+m]\times I$ implies that
\[
f^m(\mathfrak{p}\circ h_i^+) \sim g*(\mathfrak{p}\circ h_{i+m}^+)*\overline{g}'
\]
in $R$, where here
\begin{align*}
g &\sim f^m(\mathfrak{p}\circ g_{i})*f^{m-1}(\mathfrak{p}\circ g_{i+1})*\ldots *(\mathfrak{p}\circ g_{i+m})\\
g' &\sim f^m(\mathfrak{p}\circ g_i')*f^{m-1}(\mathfrak{p}\circ g_{i+1}')*\ldots *(\mathfrak{p}\circ g_{i+m}')
\end{align*}
are paths (of shortest length in their homotopy class) in $R$. We have that
\[
\len(g), \len(g') \leqslant mK^m(\rho+C).
\]
Since $\mathfrak{p}\circ h_i^+$ and $\mathfrak{p}\circ h_{i+m}^+$ are immersed paths in $R_{l}$, we may apply \cref{lem:iteration_general} with $\omega = mK^m(\rho+C)$ to conclude that 
\[
\len(h_{i+m}^+)\leqslant 2M + 1.
\]
Here we are using the fact that a geodesic in the coned-off tree of spaces connecting two vertices in a copy of $\widetilde{R}_{A_i}$ (for $1\leqslant i\leqslant n$) has length at most one. By \eqref{eq: length bound 3}, we see that
\[
\len(h\mid \{0\}\times I)\leqslant K^{2m}(2M+1 + 2(C+\rho)). 
\]
This completes the proof.
\end{proof}

\end{document}
